\definecolor{VioletD}{HTML}{534AB7}
\definecolor{VioletL}{HTML}{EEEDFE}
\definecolor{TealD}{HTML}{0A6B52}
\definecolor{TealL}{HTML}{D8F2EA}

\begin{tikzpicture}[
    font=\small,
    neuron/.style={circle, draw, thick, minimum size=0.52cm, inner sep=0pt},
    conn/.style={gray!30, thin},
]

%% ── COORDENADAS DAS CAMADAS ──────────────────────────────────────
%% i=1 no topo, cresce para baixo; espaçamento vertical 0.85cm
\foreach \i in {1,...,8}  \coordinate (IN\i)  at ( 0.0, {(4.5-\i)*0.85});
\foreach \i in {1,...,5}  \coordinate (ENC\i) at ( 3.0, {(3.0-\i)*0.85});
\foreach \i in {1,...,3}  \coordinate (LAT\i) at ( 6.0, {(2.0-\i)*0.85});
\foreach \i in {1,...,5}  \coordinate (DEC\i) at ( 9.0, {(3.0-\i)*0.85});
\foreach \i in {1,...,8}  \coordinate (OUT\i) at (12.0, {(4.5-\i)*0.85});

%% ── CONEXÕES (desenhadas antes dos nós para ficarem abaixo) ──────
\foreach \i in {1,...,8} \foreach \j in {1,...,5} \draw[conn] (IN\i)  -- (ENC\j);
\foreach \i in {1,...,5} \foreach \j in {1,...,3} \draw[conn] (ENC\i) -- (LAT\j);
\foreach \i in {1,...,3} \foreach \j in {1,...,5} \draw[conn] (LAT\i) -- (DEC\j);
\foreach \i in {1,...,5} \foreach \j in {1,...,8} \draw[conn] (DEC\i) -- (OUT\j);

%% ── NÓS (desenhados sobre as conexões) ───────────────────────────
\foreach \i in {1,...,8} {
    \node[neuron, fill=gray!20, draw=gray!55] at (IN\i)  {};
    \node[neuron, fill=gray!20, draw=gray!55] at (OUT\i) {};
}
\foreach \i in {1,...,5} {
    \node[neuron, fill=TealL, draw=TealD] at (ENC\i) {};
    \node[neuron, fill=TealL, draw=TealD] at (DEC\i) {};
}
\foreach \i in {1,...,3}
    \node[neuron, fill=VioletL, draw=VioletD] at (LAT\i) {};

%% ── RÓTULOS x_i e x̂_i ──────────────────────────────────────────
\foreach \i in {1,...,8} {
    \node[font=\scriptsize, anchor=east]  at ($(IN\i)+(-0.32,0)$) {$x_{\i}$};
    \node[font=\scriptsize, anchor=west] at ($(OUT\i)+(0.32,0)$) {$\hat{x}_{\i}$};
}

%% ── TÍTULOS (mesma altura, sem sobreposição) ─────────────────────
%% Topo do neurônio mais alto: (4.5-1)*0.85 = 2.975; folga = 0.5cm
\node[font=\small\bfseries, align=center] at ( 0.0, 3.7)
    {Camada\\de Entrada};
\node[font=\small\bfseries, align=center] at ( 3.0, 3.7)
    {Codificador};
\node[font=\small\bfseries, align=center] at ( 6.0, 3.7)
    {Representação\\Latente};
\node[font=\small\bfseries, align=center] at ( 9.0, 3.7)
    {Decodificador};
\node[font=\small\bfseries, align=center] at (12.0, 3.7)
    {Camada\\de Saída};

\end{tikzpicture}